
\section{Related Work} \label{s:relwork}

The NoSQL space has been quite active, with marked similarity among implementations. Distributed systems can be distinguished by how they choose to bias availability versus consistency. Brewer's CAP theorem \cite{Gilbert:2002:BCF:564585.564601} states that a distributed system can provide for any two of the three attributes of \emph{Consistency, Availability} and \emph{Partition tolerance}. Since networks (particularly WANs) will always partition (i.e., over time there will be some parts of the network that are temporarily unreachable), for distributed systems this distills to a choice of a design that favors strong consistency or one that is highly available in the face of partitions.

This is the heart of the distinction between SQL and NoSQL systems. SQL systems adhere to ACID properties while most NoSQL stores follow BASE properties\footnote{ACID stands for Atomicity, Consistency, Isolation and Durability; BASE stands for Basically Available, Soft state, Eventual consistency. Descriptions of ACID and BASE can be found in \cite{Wright:2007:EAS:1242520.1242521} and \cite{Cattell:2011:SSN:1978915.1978919} respectively.}. Many NoSQL implementations provide \emph{eventual consistency} \cite{Vogels:2008:EC:1466443.1466448}, where writes are permitted in the presence of partitions. In such systems a client can update any replica of an object and all updates to an object will eventually be applied, but potentially in different orders at different replicas, thus creating temporal inconsistencies that must be sorted out by client \apps{} \cite{Cooper:2008:PYH:1454159.1454167}.

Amazon's Dynamo \cite{DeCandia:2007:DAH:1323293.1294281} is the classic example of the eventually consistent model. Dynamo, a key-value store, is a zero-hop distributed hash table, where each node maintains enough routing information locally to route a request to the appropriate node directly. Dynamo shares a number of characteristics with other \dos{} used by public cloud vendors that stem from their operational model, which is a closed loop system where both clients and all services are under the control of a single company. These systems, while used to support public facing services, can assume that their \emph{internal} operating environment is non-hostile and therefore have few security requirements such as authentication and authorization. Further, since the NoSQL \oss{} permit inconsistency, users (client \apps{}) must contain an agreed understanding of how to resolve conflicts during reads. As a general-purpose system with non-specific clients, HCP can make neither assumption and must design for hostile environments and provide strong consistency guarantees.

Other systems that use  an eventual consistency model include Facebook's Cassandra (now Apache) and Google's Bigtable.  Bigtable is a distributed storage system for managing structured data. It maintains a sparse, multi-dimensional sorted map and allows applications to access their data using multiple attributes \cite{Chang:2008:BDS:1365815.1365816}. Cassandra has been described as a marriage of Dynamo and Bigtable \cite{Lakshman:2010:CDS:1773912.1773922}. Yahoo's PNUTS \cite{Cooper:2008:PYH:1454159.1454167} supports Yahoo! web properties such as Flickr. It uses per-record timeline consistency, where all replicas of a given record apply all updates to the record in the same order. Data objects (blobs) follow the eventually consistent model.

Providing strong consistency guarantees, Microsoft's Windows Azure Storage (WAS)  \cite{Huang:2012:ECW:2342821.2342823} mixes strong consistency within a \emph{local stamp} (the analog of a local cluster for HCP) with subsequent replication to remote geographies. This model is the most similar to HCP, which provides strong write guarantees through synchronous replicas locally before asynchronously dispersing replicas geographically. Such systems might be referred to as \emph{eventually distributed}.

Rather than requiring the \os{} to deal with collections of \fss{}, some implementations provide a single (logical) distributed \fs{}, such as the Google File System \cite{Ghemawat:2003:GFS:945445.945450} and CEPH \cite{Weil:2006:CSH:1298455.1298485}. By contrast, HCP uses a confederation of local file systems distributed over a collection nodes to persist data objects. This detail is abstracted from clients in favor of a storage service key-value model, where the key is a URI that indicates a unique object.
